\begin{lemma}[Left presentation invariance]
If $\mathsf{PowerPresentationEq}(X,X')$, then for every
$Y\in\mathsf{PowerQ}(Q)$,
\[
\mathsf{PowerRel}(r)(X,Y)\quad\Longleftrightarrow\quad
\mathsf{PowerRel}(r)(X',Y).
\]
\end{lemma}
\begin{proof}
We use well-founded induction on the left presentation $X$ for the child
relation \noderef{PowerChild}, which is well-founded by
\noderef{PowerChildWellFounded}.  More precisely, the induction assertion at
$X$ is that, for every $X'$ and every $Y$, a hypothesis
$\mathsf{PowerPresentationEq}(X,X')$ implies
\[
 \mathsf{PowerRel}(r)(X,Y)\longleftrightarrow
 \mathsf{PowerRel}(r)(X',Y).
\]
The induction hypothesis is therefore available whenever a child of $X$ is
compared with any structurally equivalent presentation.  The recursive
comparisons below are the ones specified by \noderef{PowerRel}; every descent
from a node $X=\mathsf{node}(\iota,h,f)$ is justified explicitly by
\[
 \mathsf{PowerChild}(f(i),X),
\]
using the witness $i$ in the definition \noderef{PowerChild}.

First suppose that $X=\mathsf{atom}(q)$.  If $X'$ is a node, then the
atom--node clause of \noderef{PowerPresentationEq} is false, contradicting
the hypothesis.  If $X'=\mathsf{atom}(q')$, that hypothesis reduces to
$q=q'$.  When $Y=\mathsf{atom}(t)$, the reductions in
\noderef{PowerRelAtomAtom} identify the two sides with $r(q,t)$ and
$r(q',t)$, respectively, and these are equal after substituting $q=q'$.
When $Y=\mathsf{node}(\lambda,k,u)$, the reduction in
\noderef{PowerRelAtomNode} identifies the two sides with
\[
 \exists \ell\in\lambda\;\mathsf{PowerRel}(r)(\mathsf{atom}(q),u(\ell))
 \quad\text{and}\quad
 \exists \ell\in\lambda\;\mathsf{PowerRel}(r)(\mathsf{atom}(q'),u(\ell)),
\]
which again agree after the same substitution.  This proves the atom case.

Now suppose that $X=\mathsf{node}(\iota,h,f)$.  A mixed $X'$ is impossible:
the node--atom clause of \noderef{PowerPresentationEq} is false.  Thus write
$X'=\mathsf{node}(\kappa,h',g)$.  Unfolding the node--node clause of
\noderef{PowerPresentationEq}, write its two matching conjuncts as
\[
 H_{\to}:\ \forall i\in\iota\;\exists j\in\kappa,
   \mathsf{PowerPresentationEq}(f(i),g(j)),
 \qquad
 H_{\leftarrow}:\ \forall j\in\kappa\;\exists i\in\iota,
   \mathsf{PowerPresentationEq}(f(i),g(j)).
\]
The equivalence properties in \noderef{PowerPresentationEqEquivalence}
permit a matching witness to be reversed by symmetry, or two such witnesses
to be composed by transitivity; the direct applications below use the
orientations displayed in $H_{\to}$ and $H_{\leftarrow}$.

Consider first $Y=\mathsf{atom}(q)$.  By the node--atom reduction
\noderef{PowerRelNodeAtom}, the desired equivalence is
\[
 \Bigl(\forall i\in\iota\;\mathsf{PowerRel}(r)(f(i),\mathsf{atom}(q))\Bigr)
 \longleftrightarrow
 \Bigl(\forall j\in\kappa\;\mathsf{PowerRel}(r)(g(j),\mathsf{atom}(q))\Bigr).
\]
For the forward implication, assume the left universal statement and fix
$j\in\kappa$.  Choose $i\in\iota$ from $H_{\leftarrow}$, and let
$e:\mathsf{PowerPresentationEq}(f(i),g(j))$ be the resulting match.  The
explicit witness $i$ gives \(\mathsf{PowerChild}(f(i),X)\) by
\noderef{PowerChild}, so the induction hypothesis at $f(i)$, applied to $e$
and the fixed right presentation $\mathsf{atom}(q)$, transfers
\(\mathsf{PowerRel}(r)(f(i),\mathsf{atom}(q))\) to
\(\mathsf{PowerRel}(r)(g(j),\mathsf{atom}(q))\).  Since $j$ was arbitrary,
the right universal statement follows.  Conversely, fix $i\in\iota$, choose
$j\in\kappa$ from $H_{\to}$, and apply the same induction hypothesis to the
explicit descent witness \(\mathsf{PowerChild}(f(i),X)\); the right
universal hypothesis supplies the comparison with $g(j)$, and the induction
hypothesis transfers it back to $f(i)$.  This proves the node--atom case
without weakening its universal left-child quantifier.

Finally let $Y=\mathsf{node}(\lambda,k,u)$.  By the node--node reduction
\noderef{PowerRelNodeNode}, the desired equivalence is
\[
 \Bigl(\forall i\in\iota\;\exists \ell\in\lambda\;
      \mathsf{PowerRel}(r)(f(i),u(\ell))\Bigr)
 \longleftrightarrow
 \Bigl(\forall j\in\kappa\;\exists \ell\in\lambda\;
      \mathsf{PowerRel}(r)(g(j),u(\ell))\Bigr).
\]
For the forward implication, fix $j\in\kappa$ and choose $i$ and
$e:\mathsf{PowerPresentationEq}(f(i),g(j))$ from $H_{\leftarrow}$.  The
assumed left statement gives a witness $\ell\in\lambda$ with
$\mathsf{PowerRel}(r)(f(i),u(\ell))$.  As above,
$\mathsf{PowerChild}(f(i),X)$ is witnessed by $i$, so the induction
hypothesis with the fixed right presentation $u(\ell)$ changes this to
$\mathsf{PowerRel}(r)(g(j),u(\ell))$.  Thus the same $\ell$ is a valid
right-child witness, and the right universal-existential statement follows.
For the reverse implication, fix $i\in\iota$ and choose $j$ and
$e:\mathsf{PowerPresentationEq}(f(i),g(j))$ from $H_{\to}$.  The assumed
right statement supplies $\ell\in\lambda$ with
$\mathsf{PowerRel}(r)(g(j),u(\ell))$.  Applying the induction hypothesis at
the explicitly witnessed child $f(i)$ transfers this comparison back to
$\mathsf{PowerRel}(r)(f(i),u(\ell))$.  This preserves the existential
right-child witness and proves the reverse implication, completing the
well-founded induction.
\end{proof}
